\scanpage{79}
plongés dans $X''$ normalisé de $X'''$ \deleted{dont} \inserted{avec lequel} on part, et \deleted{on} trouver des antiiso entre les bords \deleted{de ces} de $D'_i,D''_i$ -- d'où possibilité de recoller pour trouver une $X$ surface hol à bord, si on a aussi enlevé des disques ouv. $\mathring D_{\alpha}^{\unsure{*|k}}$ correspondants aux pts marqués $x_{\alpha}$, sinon une surface à pts marqués), découpée par des cercles $\Gamma_i$ qui ne sont \deleted{pas} \inserted{peut-être pas} néc. géodésiques \deleted{(\unsure{ce sera}} \inserted{(?? à} examiner de plus près) mais définissent des cercles géodésiques uniques, par les principes de Thurston.

Il faudrait voir que le résultat obtenu ne dépend pas des choix des paires $(D'_i,D''_i)$ mut.\textsuperscript{t} polaires, il y a un argument heuristique pour cela tout à fait vaseux, consistant à dire qu'une petite variation de \unsure{cette,|cett,}
% archive sheet number: 81
